\documentclass[12pt]{article}
\usepackage[utf8]{inputenc}
\usepackage[T1]{fontenc}
\usepackage{hyperref}
\title{Worrorra Language Resources for K-12 Curriculum Development}
\author{Compiled by Claude Sonnet 5.5 (claude-sonnet-5-5) for Andr\'as Kornai}
\date{}
\begin{document}
\maketitle

\section{Introduction}
Worrorra (ISO 639-3 wro; Glottocode worr1237; Worrorran family) was traditionally spoken on the coast and hinterland of the north-west Kimberley, Western Australia, between Walcott Inlet and the Prince Regent River. Most remaining people connected to it are at Mowanjum Community near Derby. The language may no longer have first-language speakers. The main resource is a full grammar based on the speech of the last full speakers. I found no school materials in the results.

\section{Primary Educational Resources}
I found no primers, readers or school curriculum materials specific to Worrorra. A heritage curriculum would have to be created from the grammar and archival sources below.

\section{Linguistic and Reference Materials}
\begin{itemize}
\item Mark Clendon, \emph{Worrorra: a language of the north-west Kimberley coast} (University of Adelaide Press, 2014; DOI 10.20851/worrorra; CC BY-NC-ND 4.0). It is based on the speech of Patsy Lulpunda, Amy Peters and Daisy Utemorrah, and the work was initially funded by the Kimberley Language Resource Centre. Free PDFs: \url{https://www.adelaide.edu.au/press/ua/media/709/uap-worrorra-ebook.pdf} and \url{https://library.oapen.org/bitstream/20.500.12657/33150/1/560353.pdf}.
\item The AIATSIS AUSTLANG entry lists older materials, including Love's 1920s--30s work and Coate's fieldwork. The related \emph{Handbook of Kimberley Languages} entry is at \url{https://collection.aiatsis.gov.au/austlang/language/k17}.
\end{itemize}

\section{Religious and Community Resources}
Mowanjum Arts describes the languages of the Mowanjum community (Worrorra, Ngarinyin and Wunambal) as critically endangered: \url{https://mowanjumarts.com/about/language}. I saw no Worrorra scripture and give no contact details.

\section{Digital and Audio Resources}
I found no audio or digital resources specific to Worrorra in the results.

\section{Curriculum Development Context}
Worrorra is not known to be taught in schools. Heritage and revitalisation work would rest on Clendon's grammar and its recordings of the last full speakers, plus the Mowanjum community, the Kimberley Language Resource Centre and archival sources. Related work on the neighbouring Wunambal and Gaambera languages exists, but those are different languages and are not listed here. Material for Worrorra itself is thin.
\end{document}